\section{Torque task coordinates and null-space motion}
\label{sec:task-coordinates}

The physical currents are not sacred coordinates.  For the torque surface
\(M(\vect i)=M_{\mathrm{ref}}\), the gradient
\(\vect g_M=\nabla_{\vect i}M\) is Euclidean-normal.  Every tangent velocity
satisfies \(\vect g_M^{\mathsf T}\dot{\vect i}=0\) and changes torque only at
second order.  With dynamic metric \(\matr G\), the physically fastest normal
direction is the natural gradient
\[
 \vect n_G=\frac{\matr G^{-1}\vect g_M}
 {\vect g_M^{\mathsf T}\matr G^{-1}\vect g_M}.
 \tag{29.1}
\]
It is normalized so that \(\vect g_M^{\mathsf T}\vect n_G=1\).  The associated
metric-orthogonal tangent projector is
\[
 \boxed{\gls{sym:projg}=\matr I-
 \frac{\matr G^{-1}\vect g_M\vect g_M^{\mathsf T}}
 {\vect g_M^{\mathsf T}\matr G^{-1}\vect g_M}.}
 \tag{29.2}
\]
Consequently,
\[
 \dot{\vect i}=\underbrace{\vect n_G\dot M}_{\text{torque-normal task motion}}
 +\underbrace{\matr P_G\vect v_0}_{\text{torque-tangential reallocation}}.
 \tag{29.3}
\]
The ordinary gradient asks which ampere change modifies torque most; the
natural gradient asks which dynamically cheap ampere change does so.

\begin{figure}[tb]
\centering
\begin{tikzpicture}[font=\small]
 \begin{axis}[width=.78\linewidth,height=.48\linewidth,axis lines=middle,
  xmin=-2.8,xmax=3.0,ymin=-1.7,ymax=2.3,ticks=none,xlabel={$z_\psi$},ylabel={$z_q$},clip=false]
  \addplot[Purple,very thick,domain=.35:2.7,samples=100] ({x},{1.2/x});
  \addplot[Purple,very thick,domain=-2.7:-.55,samples=100] ({x},{1.2/x});
  \addplot[only marks,mark=*,BrickRed,mark size=3pt] coordinates {(1.05,1.1429)};
  \draw[->,BrickRed,very thick] (axis cs:1.05,1.1429)--(axis cs:1.72,1.76) node[above] {$\dot{\vect i}_\perp$};
  \draw[->,ForestGreen!70!black,very thick] (axis cs:1.05,1.1429)--(axis cs:1.72,.62) node[right] {$\dot{\vect i}_\parallel$};
  \draw[blue,very thick,->] (axis cs:.55,2.18)..controls(axis cs:1.0,1.25)and(axis cs:1.65,.8)..(axis cs:2.35,.51);
  \node[Purple,anchor=south west] at (axis cs:1.9,.72) {$M=M_{\rm ref}$};
  \node[blue,anchor=west] at (axis cs:2.25,.25) {loss-improving reallocation};
 \end{axis}
\end{tikzpicture}
\caption{Normal motion changes torque; tangent motion reallocates current on a
constant-torque manifold.  The metric-weighted normal need not coincide with
the Euclidean normal drawn by the raw current axes.}
\label{fig:torque-nullspace}
\end{figure}



A task-priority controller can therefore acquire torque through the normal
component while projecting loss descent into the tangent space,
\[
 \vect v_0=-k\matr G^{-1}\nabla P_{\mathrm{Cu}}.
 \tag{29.4}
\]
This compact law is not globally time-optimal in general: drift, finite steps,
state constraints, and curvature couple the two tasks.  It is exact as the
minimum-metric-norm solution of the instantaneous equality
\(\dot M=\vect g_M^{\mathsf T}\dot{\vect i}\), and is therefore a useful local
policy and warm start.

\subsection{Hyperbolic torque and allocation rates}

The static hyperbolic coordinates become especially informative dynamically:
\[
 z_\psi=\frac{i_T}{\sqrt2}e^\alpha,\qquad
 z_q=\frac{i_T}{\sqrt2}e^{-\alpha},\qquad
 M=\frac{\kappa}{2}i_T^2.
 \tag{29.5}
\]
Differentiation gives
\[
 \dot z_\psi=z_\psi\left(\frac{\dot i_T}{i_T}+\dot\alpha\right),
 \qquad
 \dot z_q=z_q\left(\frac{\dot i_T}{i_T}-\dot\alpha\right),
 \qquad
 \dot M=\kappa i_T\dot i_T.
 \tag{29.6}
\]
Thus \(\dot i_T\) is exactly torque-changing motion, while \(\dot\alpha\)
reallocates the two torque factors at constant torque.  The third coordinate
\(z_0\) remains a second torque-null freedom.  Away from \(i_T=0\), the EESM
can be written locally as
\[
 \begin{bmatrix}\dot i_T&\dot\alpha&\dot z_0\end{bmatrix}^{\mathsf T}
 =\vect F_T+\matr G_T\vect u.
 \tag{29.7}
\]
This is the desired ``torque plus allocation'' machine.  The Jacobian becomes
singular at zero torque and can be ill-conditioned when one torque factor is
small, so a controller should switch back to regular whitened coordinates
near those sets.

\subsection{Start and target sets}

Fast torque response targets the manifold
\[
 \gls{sym:torqueset}=\{\vect i:M(\vect i)=M_{\mathrm{ref}}\},
 \tag{29.8}
\]
whereas fast convergence to efficient steady operation targets the single
point \(\vect i^\star(M_{\mathrm{ref}})\).  A hierarchical strategy reaches
\(\mathcal T\) first and then follows a torque-tangential curve toward
\(\vect i^\star\).  These terminal conditions have different minimum times;
reporting one without naming the target is ambiguous.

On \(\mathcal T\), the induced metric is simply the restriction of
\(\matr G\) to tangent vectors.  The optimal redistribution curve is a
geodesic of that induced metric only when its objective is travel time along
the manifold and drift is absent or incorporated into the correct Finsler
density.  Direct current interpolation usually leaves the surface and causes
torque error.

